\section{实验结果与分析}
\subsection{预处理热点与单项消融}
首次分项诊断显示，读取解析、编号排序去重、边排序去重和映射度数统计合计约占预处理的96\%--97\%。这些是单次诊断，不作为重复实验统计。基于此，优化集中于前三类高成本工作，而没有优先并行化占比较小的CSR扫描等步骤。
\begin{table}[htbp]
\centering
\small
\caption{预处理单项与组合消融（中位数，ms）}\label{tab:ablation}
\begin{tabular}{lrrr}
\toprule
配置 & Amazon & roadNet-CA & web-Google \\
\midrule
\texttt{legacy} & 513.175 & 2791.512 & 3168.329 \\
\texttt{fast} & 376.883 & 1984.585 & 2459.786 \\
\texttt{parse8} & 342.952 & 1696.512 & 2229.596 \\
\texttt{map1} & 520.483 & 2747.486 & 3165.870 \\
\texttt{map8} & 406.495 & 2544.171 & 2551.618 \\
\texttt{sort8} & 438.807 & 2093.300 & 2379.322 \\
\texttt{all8} & 150.166 & 779.851 & 773.687 \\
\bottomrule
\end{tabular}
\par\smallskip{\footnotesize 每配置每图3次，仅计预处理；不同于完整GPU流程主指标。配置定义见正文。}
\end{table}


快速单线程解析已获得一定收益，说明解析实现改写和多线程是两个不同因素。映射拆分单线程与原版接近，没有证据说明拆分本身稳定更快；并行映射则抵消了额外扫描成本。排序优化对道路图和网页图的收益较Amazon明显。各项加速比不能相乘，组合效果应以实际组合测量为准。

\subsection{线程数扩展}
\begin{table}[htbp]
\centering
\small
\caption{组合预处理的高线程数扫描（中位数，ms）}\label{tab:threads}
\begin{tabular}{rrrr}
\toprule
线程数 & Amazon & roadNet-CA & web-Google \\
\midrule
8 & 150.513 & 715.829 & 791.550 \\
16 & 123.352 & 608.375 & 622.032 \\
32 & 116.638 & 553.633 & 580.379 \\
64 & 118.860 & 579.019 & 575.759 \\
\bottomrule
\end{tabular}
\par\smallskip{\footnotesize 每配置每图3次；仅使用同一轮8/16/32/64线程扫描，不与早期2/4/8轮次混合。}
\end{table}

从8到32线程，三个图的预处理中位数继续下降；从32到64线程，Amazon与道路图中位数略增，Google仅有约0.8\%的降低，且三次样本的范围重叠。因此选择32线程作为统一实验配置，而不声称找到了硬件或所有图上的全局最优值。没有额外证据时，不将变化直接归因于NUMA、内存带宽或某一同步机制。

\subsection{完整任务时间}
\begin{table}[htbp]
\centering
\small
\caption{完整计数时间与优化加速比}\label{tab:full}
\begin{tabular}{lrrrr}
\toprule
数据集 & 基础版/ms & 优化版/ms & 加速比 & 耗时降低 \\
\midrule
Amazon & 933.250 & 566.650 & 1.65$\times$ & 39.3\% \\
roadNet-CA & 3255.282 & 984.945 & 3.31$\times$ & 69.7\% \\
web-Google & 3677.514 & 1038.127 & 3.54$\times$ & 71.8\% \\
\bottomrule
\end{tabular}
\par\smallskip{\footnotesize 每版每图独立完整运行20次；表中时间为中位数。优化版为32线程组合预处理，GPU边并行计数相同。}
\end{table}

在GPU计数内核相同的条件下，表~\ref{tab:full} 显示三个图的完整任务中位数时间分别约加速1.65、3.31和3.54倍。以上是本阶段主结果；预处理消融的加速比只描述局部阶段，不替代完整指标。

Amazon整体收益低于局部预处理收益，与单次记录中数百毫秒CUDA准备成本相符；但本批正式实验尚未汇总全部分项原始样本，不能仅凭总体中位数作严格成本分解。道路图和Google原始预处理更重，因此改善该部分能带来较大的完整流程收益。所有结论限于本次输入、硬件和共享运行条件，原始波动不应被静默删除。

\subsection{A100上的任务划分探索}
另在A100-PCIE-40GB上，对14图、edge/warp/split16/split32/split64五种配置各执行20次独立进程，共1400份原始记录，均通过同一CSR上的CPU双指针核验。预处理统一为12线程OpenMP；这不是前述A800预处理比较的继续样本，不跨机器混合计算加速比。
\begin{table}[htbp]
\centering
\small
\caption{A100调度对照：GPU计数阶段}\label{tab:a100_gpu}
\begin{tabular}{lrrrrr}
\toprule
图 & edge/ms & warp比值 & split16比值 & split32比值 & split64比值 \\
\midrule
C1 & 0.1012 & 0.230 & 0.883 & 0.849 & 0.891 \\
C2 & 0.1173 & 0.109 & 0.870 & 0.904 & 0.838 \\
C3 & 0.9735 & 0.347 & 1.032 & 1.011 & 1.113 \\
W1 & 0.0959 & 0.212 & 0.865 & 0.826 & 0.850 \\
W2 & 0.1688 & 1.111 & 0.932 & 0.931 & 0.877 \\
W3 & 0.4317 & 0.883 & 0.995 & 0.916 & 0.995 \\
W4 & 4.5490 & 0.319 & 0.978 & 0.988 & 0.999 \\
S1 & 0.1361 & 0.585 & 0.818 & 0.883 & 0.859 \\
S2 & 0.0963 & 0.083 & 0.862 & 0.850 & 0.853 \\
S3 & 0.0595 & 0.243 & 0.776 & 0.771 & 0.750 \\
S4 & 1.6754 & 1.052 & 1.026 & 0.999 & 0.987 \\
S5 & 22.7226 & 1.002 & 0.853 & 0.924 & 0.956 \\
S6 & 0.8289 & 0.711 & 1.371 & 1.258 & 1.258 \\
S7 & 3.5839 & 1.102 & 1.190 & 1.077 & 1.119 \\
\bottomrule
\end{tabular}
\par\smallskip{\footnotesize 每图每配置20次独立进程；时间为中位数，比值为edge中位数除以对应配置中位数。大于1为更快。Wiki-Vote默认blocks不同。}
\end{table}

固定warp在14图中的10图GPU计数阶段更慢；split16在S6、S7上分别约1.37、1.19倍加速，而S5上退化。结果支持收益依赖任务结构，不支持分流普遍优于边并行。队列和协作的成本尚需分项诊断，不能仅凭总时间断定队列是唯一瓶颈。
\begin{table}[htbp]
\centering
\small
\caption{A100调度对照：完整任务}\label{tab:a100_full}
\begin{tabular}{lrrrrr}
\toprule
图 & edge/ms & warp比值 & split16比值 & split32比值 & split64比值 \\
\midrule
C1 & 413.345 & 1.007 & 1.005 & 1.013 & 0.998 \\
C2 & 994.656 & 1.036 & 1.022 & 1.020 & 1.015 \\
C3 & 1087.742 & 1.019 & 1.022 & 1.018 & 1.036 \\
W1 & 433.727 & 0.983 & 0.993 & 0.959 & 0.985 \\
W2 & 315.392 & 1.056 & 1.002 & 1.019 & 1.041 \\
W3 & 359.332 & 1.038 & 1.052 & 1.057 & 1.055 \\
W4 & 3362.460 & 1.014 & 1.019 & 0.991 & 0.999 \\
S1 & 319.795 & 1.010 & 1.014 & 1.008 & 0.957 \\
S2 & 535.588 & 1.034 & 1.016 & 1.027 & 1.016 \\
S3 & 360.926 & 1.028 & 1.029 & 0.989 & 0.952 \\
S4 & 405.383 & 1.040 & 0.994 & 1.036 & 1.026 \\
S5 & 778.532 & 0.999 & 0.983 & 0.997 & 0.989 \\
S6 & 535.273 & 1.031 & 1.053 & 1.041 & 1.012 \\
S7 & 529.955 & 1.047 & 1.039 & 1.020 & 1.046 \\
\bottomrule
\end{tabular}
\par\smallskip{\footnotesize 每图每配置20次独立进程；时间为中位数，比值为edge中位数除以对应配置中位数。大于1为更快。Wiki-Vote默认blocks不同。}
\end{table}

完整任务包含预处理、CUDA准备、传输和全部计数归约，各配置的小幅差异不能全部归因GPU调度。280个edge样本的CUDA准备时间整体中位数约295.74ms。Wiki-Vote的edge默认394块、其他方案864块，需补充相同网格对照；其他输入均为864块，每块256线程。

原始包及重算脚本见仓库\texttt{results/a100\_v1/}。图数据与可执行文件不在包内，输入与二进制哈希仅为运行时记录。CPU型号、频率和硬件独占性未完整建立；本轮仍属于探索性评测，不宣称获得全局最优参数。
